\documentclass[12pt,a4paper]{article}
\usepackage[utf8]{inputenc}
\usepackage[T2A]{fontenc}
\usepackage[russian]{babel}
\usepackage{amsmath,amssymb,amsthm}
\usepackage{booktabs,array}
\usepackage{microtype}
\usepackage{geometry}
\geometry{margin=2.3cm}
\setlength{\emergencystretch}{2em}

\newtheorem{proposition}{Предложение}
\newtheorem{nogocriterion}{Критерий провала}

\title{Spin-структуры и счёт поколений\\
Третий gate гипотезы меню состояний}
\author{}
\date{4 августа 2026 г.}

\begin{document}
\maketitle

\section{Проблема}

После положительного \(SU(5)\)-gate остаётся вопрос о трёх поколениях. Для
\[
K=\mathbb{RP}^3\times S^1
\]
пространство spin-структур является torsor над
\[
H^1(K,\mathbb Z_2)\cong\mathbb F_2^2,
\]
поэтому содержит четыре элемента. Формула \(4-1=3\) приобретает смысл только если один reference sector выделяется до сравнения с наблюдаемым числом поколений.

\section{Bounding reference}

Пространство \(\mathbb{RP}^3=L(2,1)\) является границей ориентированного disk bundle над \(S^2\) с Euler number \(2\). Это заполнение стягивается на \(S^2\), имеет тривиальную первую cohomology и является spin, поскольку Euler class чётен. Следовательно, его единственная spin-структура выделяет одну из двух spin-структур границы.

На окружности диск \(D^2\) аналогично выделяет bounding, или antiperiodic, spin-структуру. Поэтому product rule выделяет ровно одну пару
\[
s_{ref}=(s_{\mathbb{RP}^3}^{bound},s_{S^1}^{bound})
\]
среди четырёх product spin-структур.

\begin{proposition}[Геометрический счёт трёх nonreference sectors]
Если стандартные factor fillings являются частью menu-модели, то существует одна factorwise bounding reference spin-структура и ровно три nonreference spin-сектора.
\end{proposition}

Это первый случай, когда число три возникает не простым вычитанием выбранного элемента, а после предварительного геометрического критерия.

\section{Абстрактная симметрия тройки}

После выбора reference torsor превращается в vector space \(\mathbb F_2^2\). Его три ненулевых элемента переставляются группой
\[
GL(2,2),
\qquad
|GL(2,2)|=6,
\qquad
GL(2,2)\cong S_3.
\]
Таким образом, на абстрактном уровне три nonreference labels образуют естественную orbit и допускают permutation symmetry трёх поколений.

\section{Геометрический obstruction}

Но фундаментальная группа до редукции по модулю два равна
\[
\pi_1(K)=\mathbb Z_2\times\mathbb Z.
\]
Её автоморфизм обязан сохранять единственный torsion subgroup \(\mathbb Z_2\). Поэтому не вся группа \(GL(2,2)\) поднимается до геометрической автоморфизмной группы \(K\).

На spin labels \((p,q)\in\mathbb F_2^2\) допустимое нетривиальное преобразование имеет вид
\[
(p,q)\longmapsto(p,q+p).
\]
Три ненулевых labels распадаются на орбиты размеров
\[
1+2,
\]
а не на одну orbit размера три. Чистый circle twist выделяется отдельно, тогда как два сектора с \(\mathbb{RP}^3\)-torsion twist образуют пару.

\begin{nogocriterion}[Провал автоматической эквивалентности поколений]
Топология \(K\) даёт правильное число трёх nonreference spin-секторов, но не делает их тремя эквивалентными поколениями. Полная \(S_3\)-симметрия требует дополнительного emergent действия, не следующего из \(\operatorname{Aut}(\mathbb Z_2\times\mathbb Z)\).
\end{nogocriterion}

\section{Спектральная проверка различимости}

Две spin-структуры круглого \(\mathbb{RP}^3\) имеют противоположные Dirac \(\eta\)-invariants
\[
\eta_D=\pm\frac14.
\]
На \(S^1\) bounding структура имеет half-integer gap, а nonbounding periodic структура содержит нулевую моду. Следовательно, четыре product structures спектрально различимы; их вырождение в поколенческую тройку не возникает автоматически.

\section{Вердикт}

Третий gate даёт смешанный, но сильный результат:
\begin{itemize}
    \item уникальная bounding reference structure может быть выведена из стандартных заполнений;
    \item после её удаления остаются ровно три spin-sector labels;
    \item абстрактно они несут естественное действие \(S_3\);
    \item геометрические автоморфизмы самого \(K\) реализуют только orbit split \(1+2\).
\end{itemize}

Следующий gate должен проверить, может ли уже найденный \(SU(5)\)-fiber и его algebra переходов восстановить полное \(S_3\)-действие на трёх nonreference sectors без отдельного постулата family symmetry. Если нет, идея трёх поколений закрывается на уровне красивого счёта, но не физического механизма.

\end{document}